\documentclass[11pt,a4paper]{article}
\usepackage[utf8]{inputenc}
\usepackage[T1]{fontenc}
\usepackage[french]{babel}
\usepackage[margin=1.7cm]{geometry}
\usepackage{graphicx}
\usepackage{xcolor}
\usepackage{mdframed}
\definecolor{NavyB}{HTML}{123A75}
\definecolor{AccR}{HTML}{A81E12}
\definecolor{GrnD}{HTML}{0C5730}
\pagestyle{empty}
\setlength{\parindent}{0pt}
\newmdenv[linecolor=NavyB!35,backgroundcolor=NavyB!4,linewidth=0.8pt,
  roundcorner=3pt,innermargin=6pt,innertopmargin=5pt,innerbottommargin=5pt,
  skipabove=2mm,skipbelow=5mm]{expl}
\newcommand{\fiche}[4]{%
\begin{center}\setlength{\fboxsep}{0pt}\setlength{\fboxrule}{0.5pt}%
\fbox{\includegraphics[width=0.72\linewidth]{../images/#1.png}}\end{center}
\vspace{-1mm}
\begin{expl}
{\color{NavyB}\bfseries #2}\par\vspace{0.8mm}
{\bfseries\color{AccR}Ce qu'elle montre.} #3\par\vspace{0.8mm}
{\bfseries\color{GrnD}Quand la sortir.} #4
\end{expl}}
\begin{document}
\begin{center}
{\color{NavyB}\Large\bfseries Les 31 fiches écran expliquées\par}
\vspace{1mm}
{\small Mécathermo Ma0 --- Casio fx-CG50\\[1mm]
Chaque fiche est reproduite à sa taille réelle ($384\times216$ px), suivie de ce qu'elle contient
et du moment où il faut l'ouvrir.\\
\textbf{D01--D23} = démarche (l'ordre dans lequel rédiger).
\textbf{S01--S08} = schéma (le dessin à avoir en tête).}
\end{center}
\vspace{3mm}
{\color{NavyB}\large\bfseries Fiches DÉMARCHE}\par\vspace{2mm}
\fiche{D01-essence-volumes-points}{Moteur essence : volumes et premiers points}{Les 4 premieres lignes du tableau d etat d un moteur a essence : les volumes a partir de la cylindree et de $\varepsilon$, la quantite de matiere $n$, puis le point C par la compression adiabatique.}{Des que l enonce parle d un moteur essence avec une cylindree et un rapport volumetrique. C est toujours par la que l on commence.}
\fiche{D02-essence-combustion-rendement}{Moteur essence : fin du cycle}{La suite immediate de D01 : une fois $Q_{CD}$ connu, on obtient le point D (combustion isochore), le point E (detente) et le rendement theorique.}{Apres avoir calcule $Q_{CD}$ avec les fiches D06 et D07. Enchainer directement sur D20 pour les rendements reels.}
\clearpage
\fiche{D03-diesel-ce-qui-change}{Diesel : les 4 differences}{Ce qui distingue le diesel de l essence, et rien d autre : $\varepsilon$ plus grand, air seul admis, combustion isobare, et le nouveau parametre $\rho$.}{A lire AVANT de commencer un exercice diesel. Elle evite l erreur la plus couteuse : utiliser $C_v$ pour la combustion alors qu elle est isobare.}
\fiche{D04-diesel-points-rendement}{Diesel : points et rendement}{Les points D et E du cycle diesel, la chaleur de combustion avec $C_p$, et la formule de rendement (celle avec $\rho$).}{Juste apres D03. Le calcul des volumes, de $n$ et du point C se fait avec D01, il est identique a l essence.}
\clearpage
\fiche{D05-moteur-2-temps}{Moteur 2 temps}{Rappelle que le cycle thermodynamique est le meme que l essence, et que seule la formule de puissance change.}{Des que l enonce dit 2 temps. C est la fiche anti-piege : elle vous empeche de garder le $/2$ du 4 temps.}
\fiche{D06-combustion-reaction}{Combustion : equilibrer et passer a l air}{L equation de combustion d un hydrocarbure, puis le passage des moles d oxygene aux moles d air (l air ne contient que 21 pour cent d O2), et l application de l exces d air $\lambda$.}{Des que l enonce donne une formule de carburant du type CH$_{1,8}$ et un $\lambda$.}
\clearpage
\fiche{D07-combustion-chaleur}{Combustion : jusqu a la chaleur}{La suite : repartir les moles du cylindre entre air et carburant, convertir en masse, puis obtenir $Q_{CD}$ avec le PCI.}{Immediatement apres D06. Attention a la derniere ligne : le PCI est donne en kJ/kg, il faut le multiplier par 1000.}
\fiche{D08-exces-air-consequences}{Exces d air : les consequences}{Ce qui se passe quand $\lambda$ est trop petit (CO puis imbrules) ou trop grand (air comprime pour rien).}{C est une question de cours, pas un calcul. A relire avant l examen si une question theorique sur $\lambda$ tombe.}
\clearpage
\fiche{D09-frigo-placer-points}{Frigo : placer les 4 points}{La marche a suivre pour placer 1', 2', 3' et 4' sur le diagramme des frigoristes, avec la surchauffe et le sous-refroidissement.}{Des qu on vous donne un diagramme (log p, h). C est la fiche la plus utile de toutes : sans elle, impossible de lire les enthalpies.}
\fiche{D10-frigo-energies}{Frigo : les trois energies}{Correction du point 2 avec le rendement isentropique, puis $w$, $q_1$ et $q_2$ comme simples differences d enthalpie.}{Une fois les 4 points places. La derniere ligne donne la verification $q_1+q_2+w=0$ : faites-la systematiquement.}
\clearpage
\fiche{D11-psn-cop}{PSN ou COP ?}{Le choix de la formule selon que la machine est un frigo (utile = chaleur prise au froid) ou une pompe a chaleur (utile = chaleur cedee au chaud).}{Au moment de repondre a la question de performance. C est ici qu on perd des points en recitant sans reflechir a ce qui est utile.}
\fiche{D12-carnot-debits}{Carnot et debits}{Les bornes theoriques $PSN_{max}$ et $COP_{max}$, puis le passage des kJ/kg aux kW via le debit.}{En fin d exercice frigo : d abord pour verifier que votre COP est credible, ensuite des que l enonce donne une puissance en kW.}
\clearpage
\fiche{D13-turbine-joule}{Turbine a gaz : les regles}{Les deux seules formules utiles en machine ouverte : adiabatique (rapport de pression, $w=c_p\Delta T$) et isobare ($q=c_p\Delta T$).}{Des que l exercice parle de compresseur, turbine et debit massique. On raisonne par kg, sans moles ni volumes.}
\fiche{D14-turbine-astuce}{Turbine : l astuce du sujet}{La traduction de la phrase piege << tout le travail de la turbine sert au compresseur >> : elle donne $T_3-T_4=T_2-T_1$.}{Quand une turbine est liee au compresseur. Rappelle aussi que ce travail ne compte pas dans le travail net.}
\clearpage
\fiche{D15-transfo-isobare-isochore}{Transformation a $p$ ou $V$ constant}{$W$ et $Q$ pour une isobare et pour une isochore, avec la relation entre les etats.}{Quand on demande le travail et la chaleur d une seule transformation, ou pour remplir une ligne du tableau d energies d un cycle.}
\fiche{D16-transfo-isotherme-adiab}{Transformation a $T$ constant ou $Q=0$}{$W$ et $Q$ pour une isotherme et pour une adiabatique reversible, avec les trois ecritures equivalentes de l adiabatique.}{Meme usage que D15. Les trois formes de l adiabatique servent selon les variables connues dans l enonce.}
\clearpage
\fiche{D17-constantes}{Les constantes a poser}{$C_v$, $C_p$, la relation de Mayer, et les valeurs numeriques pour $\gamma=1,4$.}{Tout au debut de l exercice, des que $\gamma$ est donne. On les reutilise a chaque ligne : autant les ecrire une fois pour toutes.}
\fiche{D18-entropie}{Variation d entropie}{Les quatre formules de $\Delta S$ selon le type de transformation, et le fait que $\Delta S=0$ sur un cycle.}{Si une question demande explicitement une variation d entropie. Rappelle aussi pourquoi une adiabatique reversible est dite isentropique.}
\clearpage
\fiche{D19-echange-thermique}{Echange thermique}{$Q=mc\Delta T$, la chaleur latente, et surtout $P=\dot m c \Delta T$ pour un circuit a debit.}{Pour la partie << circuit d eau >> ou << air d un local >> qui accompagne souvent l exercice frigo. Le lien entre les deux circuits est la puissance.}
\fiche{D20-rendements-cascade}{Cascade des rendements}{L enchainement $\eta_{th}$ puis $\eta_{forme}$ puis $\eta_{meca}$, et le rendement effectif global.}{Des que l enonce donne un rendement de forme et un rendement mecanique. On applique les pertes l une apres l autre, dans cet ordre.}
\clearpage
\fiche{D21-puissance}{Du travail a la puissance}{Les deux formules de puissance (4 temps et 2 temps) et la formule inverse pour retrouver le regime.}{A la derniere question d un exercice moteur. Verifiez toujours si $W$ est bien par cycle ET par cylindre avant de multiplier par $z$.}
\fiche{D22-verifications}{Les verifications avant de rendre}{Les cinq controles qui rattrapent une erreur, dont le plus puissant : la somme des travaux et des chaleurs doit etre nulle sur un cycle.}{Systematiquement, avant de passer a l exercice suivant. Chaque controle prend dix secondes.}
\clearpage
\fiche{D23-ordres-de-grandeur}{Ordres de grandeur}{Les valeurs plausibles pour une temperature de combustion, une pression maximale, un rendement, un COP et un regime moteur.}{Quand un resultat vous parait bizarre. Si vous trouvez un COP de 40 ou un rendement de 95 pour cent, il y a une erreur.}
\clearpage
{\color{NavyB}\large\bfseries Fiches SCHÉMA}\par\vspace{2mm}
\fiche{S01-cycle-essence}{Cycle essence en diagramme $p$-$V$}{Le cycle de Beau de Rochas avec le code couleur : rouge pour les adiabatiques, vert pour les isochores, bleu pour l isobare.}{Quand on demande de tracer le cycle, et comme aide-memoire pour savoir quelle transformation relie deux points.}
\fiche{S02-cycle-diesel}{Cycle diesel en diagramme $p$-$V$}{Le meme cycle, mais avec le palier isobare C-D (en bleu) au lieu de la verticale isochore de l essence.}{Meme usage. La comparaison visuelle avec S01 est le meilleur moyen de retenir la difference essence / diesel.}
\clearpage
\fiche{S03-quatre-transformations}{Les 4 courbes par un meme point}{Isochore verticale, isobare horizontale, isotherme et adiabatique en hyperboles, l adiabatique etant toujours plus pentue.}{C est une question de l examen telle quelle (question 22 des questions preliminaires). A savoir redessiner de tete.}
\fiche{S04-schema-frigo}{Les organes d une machine frigo}{Les 4 elements en boucle fermee : evaporateur, compresseur, condenseur, detendeur, avec la numerotation 1 a 4.}{C est aussi une question de l examen telle quelle (question 34). Le sens du parcours donne l ordre des transformations.}
\clearpage
\fiche{S05-diagramme-logp-h}{Le cycle sur le diagramme des frigoristes}{La forme du cycle une fois place : compression qui monte a droite, condensation horizontale, detente verticale, evaporation horizontale.}{A garder sous les yeux pendant la lecture du diagramme reel, pour verifier que votre trace a la bonne allure.}
\fiche{S06-moteur-ou-recepteur}{Moteur ou recepteur ?}{Le sens de parcours qui distingue un cycle moteur (horaire, $W<0$) d un cycle recepteur (antihoraire, $W>0$).}{Question de l examen (question 17). Sert aussi de controle de signe a la fin de n importe quel exercice de cycle.}
\clearpage
\fiche{S07-ferme-ou-ouvert}{Systeme ferme ou ouvert}{La distinction visuelle entre un gaz enferme dans un cylindre et un fluide qui traverse une machine, avec la formule associee.}{Chaque fois que vous hesitez entre $C_v$ et $C_p$. C est l erreur numero un de toute la matiere.}
\fiche{S08-cycle-joule}{Cycle de Joule en diagramme $p$-$V$}{Les deux adiabatiques et les deux isobares du cycle d une turbine a gaz.}{Pour situer les points d un exercice de turbine, et pour ne pas confondre ce cycle avec celui d un moteur a piston.}
\end{document}
